%%=============================================================================
%% Appendix
%%=============================================================================
\chapter{Appendix}
\label{ch:appendix}

%todo: move many code fragments here (Ollama section)

%=============================================================================
% Ollama Section todo: move the codefragments from Ollama and the outputs here
% I should move each test for each subchapter in Ollama -> code + 'vraag'+'output' + output here (references are already set to work after move and semantically correct)
%=============================================================================

\section{Ollama Lokaal (offline) resultaten}

\begin{listing}
    \begin{minted}{python}
        ####### Prep #######
        # als er meerdere python versies zijn: zet python interpreter op niewste(waar pip installs naar wijzen) (3.12.9)
        # pip install ollama
        # kies een model op: https://ollama.com/search
        # download de gekozen ollama model (bvb qwen3-vl:8b-instruct) via deze shell commando: "ollama run qwen3-vl:8b-instruct"
        # info: je kan ook modellen downloaden via de GUI van de ollama app, maar dan heb je geen enkele feedback of kennis van wanneer het model
        # volledig is gedownload
        # nu kan je deze file runnen (ook offline)
        
        import time
        
        t1 = time.time()
        
        ################################# code verkregen van: https://github.com/ollama/ollama-python
        
        from ollama import chat
        from ollama import ChatResponse
        
        response: ChatResponse = chat(model='qwen3-vl:8b-instruct', messages=[
        {
            'role': 'user',
            'content': 'Why is the sky blue?',
        },
        ])
        print(response['message']['content'])
        # or access fields directly from the response object
        # print(response.message.content)
        
        #################################
        
        t2 = time.time()
        
        print("\nResponse tijd in seconden: %.5f " % (t2 - t1))
    \end{minted}
    \caption[Ollama Lokaal Code]{Code om te communiceren met \textit{qwen3-vl:8b-instruct}.}
    \label{lst:ollama-lokaal-test}
\end{listing}

\textbf{Vraag:} 'Why is the sky blue?'

\textbf{Output:} 

\begin{minted}{markdown}
    The sky appears blue due to a phenomenon called **Rayleigh scattering**. Here's how it works:
    
    ### 1. **Sunlight is White, But Made of Many Colors**
    Sunlight, or white light, is actually composed of all the colors of the visible spectrum—red, orange, yellow, green, blue, indigo, and violet. These colors travel at different wavelengths.
    
    ### 2. **Atmosphere Scatters Shorter Wavelengths More**
    The atmosphere contains molecules (like nitrogen and oxygen) and tiny particles. When sunlight enters the atmosphere, these molecules scatter the sunlight in all directions.
    
    - **Shorter wavelengths** (like **blue and violet**) are scattered much more effectively than longer wavelengths (like red and orange).
    - This is because the scattering intensity is **inversely proportional to the fourth power of the wavelength** (Rayleigh scattering law: *I ∝ 1/λ⁴*).
    
    ### 3. **Why We See Blue, Not Violet**
    Although **violet** light is scattered even more than blue, our eyes are **less sensitive to violet**, and some of the violet light is absorbed by the upper atmosphere. So the dominant color we perceive is **blue**.
    
    ### 4. **What Happens at Sunrise or Sunset?**
    At sunrise or sunset, sunlight travels through more atmosphere. This causes **most of the blue and violet to be scattered out of our line of sight**, so the remaining light (with longer wavelengths like red and orange) reaches us — that’s why sunsets are orange or red.
    
    ---
    
    ✅ In summary:
    **The sky is blue because molecules in the atmosphere scatter shorter-wavelength blue light more than other colors, and our eyes are most sensitive to blue — making it the color we see most often in the daytime sky.**
    Response tijd in seconden: 84.42009
\end{minted}

\section{Ollama Cloud-Signin (online) resultaten}

\begin{listing}
    \begin{minted}{python}
        # setup:
        # run following 2 commands in terminal:
        
        # 1)
        # ollama signin
        # generated the following for me:
        # https://ollama.com/connect?name=LAPTOP-1CQ3QTCS&key=c3NoLWVkMjU1MTkgQUFBQUMzTnphQzFsWkRJMU5URTVBQUFBSUorK245dlBpY0JOVzZqZnFWRDF1Wm1Ldnh4eitYc0RwN2RkQzRZb2tJSC8
        # (follow the link and connect this device with your ollama account)
        
        # 2)
        # ollama pull gpt-oss:120b-cloud
        # save metadata about the cloud model on your system (this does not install the 65gb model locally, only some metadata)
        # "manifest" sha256
        
        # https://github.com/ollama/ollama-python?tab=readme-ov-file#cloud-models
        
        import time
        t1 = time.time()
        
        from ollama import Client
        
        client = Client()
        
        messages = [
        {
            'role': 'user',
            'content': 'Why is the sky blue?',
        },
        ]
        
        for part in client.chat('gpt-oss:120b-cloud', messages=messages, stream=True):
        print(part.message.content, end='', flush=True)
        
        t2 = time.time()
        print("\nResponse tijd in seconden: %.5f " % (t2 - t1))
    \end{minted}
    \caption[Ollama Cloud-Signin Code]{Code om te communiceren met \textit{gpt-oss:120b-cloud}.}
    \label{lst:ollama-cloud-signin-test}
\end{listing}

\textbf{Vraag:} 'Why is the sky blue?'

\textbf{Output:} 

\begin{minted}{markdown}
    **Short answer:**  
    The sky appears blue because molecules and tiny particles in Earth’s atmosphere scatter short‑wavelength (blue) light from the Sun far more efficiently than they scatter longer‑wavelength (red) light. Our eyes see that scattered blue light coming from every direction overhead.
    
    ---
    
    ## The physics behind it
    
    | Step | What happens | Why it matters |
    |------|--------------|----------------|
    | **1. Sunlight is white** | Sunlight contains a continuous spectrum of colors—from violet (~380 nm) to red (~700 nm). | “White” light is just a mix of all visible wavelengths. |
    | **2. Light enters the atmosphere** | As sunlight passes through the lower ~100 km of the atmosphere, it encounters nitrogen (N₂), oxygen (O₂), and trace gases. | These molecules are far smaller than the wavelength of visible light. |
    | **3. Rayleigh scattering** | Light is scattered by particles that are much smaller than its wavelength. The scattering intensity follows **\(I \propto \frac{1}{\lambda^4}\)** (the Rayleigh law). | Shorter wavelengths (blue, violet) get scattered 10‑16 times more strongly than red. |
    | **4. Human eye sensitivity** | The eye’s photoreceptor (cones) are most responsive to wavelengths around 555 nm (green) and less to violet (< 400 nm). The eye also has three cone types (S, M, L) that together weight blue (≈450 nm) more than violet. | Even though violet is scattered the most, we perceive the sky as blue because our eyes are less sensitive to violet and the atmosphere absorbs much of it. |
    | **5. Observed sky color** | The scattered blue light reaches us from every direction, giving the dome overhead its characteristic hue. | The direct Sun’s beam stays relatively unscattered, so we see the Sun as white/yellow while the background sky is blue. |
    
    ---
    
    ## Why the sky isn’t violet (even though violet scatters more)
    
    1. **Eye sensitivity:** Our photoreceptors are ~5–10× less sensitive to violet than to blue.  
    2. **Solar spectrum:** The Sun emits less energy at violet wavelengths than at blue.  
    3. **Atmospheric absorption:** Ozone and other gases absorb part of the violet/ultraviolet light before it can be scattered back to us.
    
    Combined, these effects tip the balance toward the **blue** we perceive.
    
    ---
    
    ## When the sky changes color
    
    | Situation | What changes | Resulting color |
    |-----------|-------------|-----------------|
    | **Sun low on the horizon (sunrise / sunset)** | Light traverses a much longer atmospheric path → more scattering of blue & green, leaving reds & oranges. | Warm reds, pinks, and oranges dominate. |
    | **Dust, pollution, or volcanic ash** | Larger particles cause **Mie scattering**, which is less wavelength‑dependent. | Sky can look whitish, hazy, or take on a brownish tint. |
    | **Clear, dry air at high altitude** | Fewer scattering molecules, less water vapor → less overall scattering. | Sky can appear a deeper, almost **ultramarine** blue. |
    | **Night** | No direct sunlight; only scattered starlight and occasional moonlight. | Sky is dark; only faint airglow or auroral emissions may be visible. |
    
    ---
    
    ## A quick mental picture
    
    1. Imagine shining a flashlight (white light) into a foggy room. The fog’s tiny water droplets scatter all colors more or less equally → the room glows white.  
    2. Now replace the droplets with **tiny gas molecules** that are orders of magnitude smaller. They act like a “blue‑biased” fog, sending far more blue photons toward your eyes from every direction. That’s the sky we see on a clear day.
    
    ---
    
    ### TL;DR
    
    - Sunlight = all colors.  
    - Molecules in the atmosphere scatter light **more efficiently at short wavelengths** (Rayleigh scattering ∝ 1/λ⁴).  
    - Blue light (≈450 nm) is scattered ~10× more than red (≈650 nm).  
    - Our eyes are tuned to see blue better than violet, and violet is partly absorbed.  
    - The net effect: a sky bathed in scattered blue light.  
    
    That’s why, on a clear day, you look up and see a beautiful blue dome. 🌤️
    Response tijd in seconden: 6.93436 
\end{minted}

\section{Ollama Cloud-API (online) resultaten}

\begin{listing}
    \begin{minted}{python}
        OLLAMA_API_KEY="d22e2e0258e749dd82b3db75a57c0fea.TNwiekYbQv1k1P_mv-vUTguW"
        
        import time
        t1 = time.time()
        
        '''
        normale procedure(voor security wil je een geheim api-key):
        
        export OLLAMA_API_KEY=your_api_key
        
        import os
        
        os.environ.get('OLLAMA_API_KEY')
        '''
        
        # https://github.com/ollama/ollama-python?tab=readme-ov-file#cloud-api-ollamacom
        from ollama import Client
        
        client = Client(
        host='https://ollama.com',
        headers={'Authorization': 'Bearer ' + OLLAMA_API_KEY}
        )
        
        messages = [
        {
            'role': 'user',
            'content': 'Why is the sky blue?',
        },
        ]
        # model: 'gpt-oss:120b' also works, not sure why, but -cloud suffix not needed
        for part in client.chat('gpt-oss:120b-cloud', messages=messages, stream=True):
        print(part.message.content, end='', flush=True)
        
        t2 = time.time()
        print("\nResponse tijd in seconden: %.5f " % (t2 - t1))
    \end{minted}
    \caption[Ollama Cloud API Code]{Code om te communiceren met \textit{gpt-oss:120b-cloud}.}
    \label{lst:ollama-cloud-api-test}
\end{listing}

\textbf{Vraag:} 'Why is the sky blue?'

\textbf{Output:}

\begin{minted}{markdown}
    **Short answer:**  
    The sky looks blue because molecules and tiny particles in Earth’s atmosphere scatter short‑wavelength (blue) light from the Sun much more efficiently than they scatter the longer‑wavelength (red, orange, yellow) colors. Our eyes receive a lot of that scattered blue light from every direction, so the whole dome overhead appears blue.
    
    ---
    
    ## The physics behind the color
    
    | Step | What happens | Why it matters for color |
    |------|--------------|---------------------------|
    | 1. Sunlight reaches Earth | Sunlight is a mixture of all visible colors (roughly a black‑body spectrum at ~5,800 K). | The light contains roughly equal amounts of red, green, blue, etc. |
    | 2. Light enters the atmosphere | Photons encounter air molecules (N₂, O₂) and very small particles (dust, water droplets). | These scatter the photons away from the straight‑line path. |
    | 3. **Rayleigh scattering** dominates | For particles **much smaller** than the wavelength of light (≈ 0.1 µm vs. 0.4–0.7 µm), the scattering cross‑section follows \(\sigma \propto \lambda^{-4}\). | Short‑wavelength (blue, ~450 nm) light is scattered **≈ 10 times** more than red (~650 nm). |
    | 4. Scattered light reaches our eyes | The scattered photons come to us from every part of the sky, not just the Sun’s direction. | Because blue photons are far more numerous after scattering, the sky’s overall hue is blue. |
    | 5. Direct sunlight stays relatively unscattered | Most of the Sun’s beam goes straight through the atmosphere; only a tiny fraction is removed. | The Sun itself appears yellowish/white because the scattering removes a bit more blue than red. |
    
    ### Rayleigh’s \(\lambda^{-4}\) law (very basic)
    
    \[
    I_{\text{scattered}}(\lambda) \;=\; I_0(\lambda)\,\frac{2\pi^5}{3}\,\frac{(n^2-1)^2}{N^2\lambda^4}\,L
    \]
    
    * \(I_0(\lambda)\) – incident intensity at wavelength \(\lambda\)  
    * \(n\) – refractive index of air  
    * \(N\) – molecular density  
    * \(L\) – path length through the scattering medium  
    
    The \(\lambda^{-4}\) term makes the intensity drop off sharply with increasing wavelength, giving blue its advantage.
    
    ---
    
    ## Why isn’t the sky black?
    
    If there were no atmosphere, we would see the blackness of space (no scattering). The atmosphere *fills* the line of sight with light that’s been redirected from the Sun. Even though each individual molecule scatters only a tiny fraction of the incident photons, the **enormous number** of molecules (≈ 2.5 × 10²⁵ per cubic meter at sea level) adds up to a noticeable glow.
    
    ---
    
    ## Variations in sky color
    
    | Condition | What changes | Resulting color |
    |-----------|--------------|-----------------|
    | **Sunrise / Sunset** | Light travels through **more atmosphere** → longer path → more scattering of blue & green → only the longest wavelengths (red–orange) survive. | Sky and Sun appear red/orange. |
    | **High humidity / haze** | Small water droplets (Mie scattering) are **size‑comparable** to visible wavelengths; scattering becomes less wavelength‑selective. | Sky looks whiter or milky. |
    | **Air pollution (smog, dust)** | Larger particles increase Mie scattering; sometimes preferentially scatter red light (e.g., after volcanic eruptions). | Sky can adopt pink, orange, or even brown tones. |
    | **Polar regions** | Low Sun angle + very clean, cold air can produce vivid indigo or violet hues due to even stronger scattering of short wavelengths. | Deep blue to violet skies. |
    | **Night** | Sun is below the horizon → no direct daylight to scatter → sky is black except for scattered moonlight and airglow. | Dark night sky. |
    
    ---
    
    ## A quick “back‑of‑the‑envelope” estimate
    
    Take a column of air 1 km thick at sea level. The probability that a given photon is scattered at least once in that column is roughly:
    
    \[
    P \approx 1 - e^{- \tau},\quad \tau = n\,\sigma L
    \]
    
    * \(n \approx 2.5\times10^{25}\,\text{m}^{-3}\) (molecules)  
    * \(\sigma_{\text{blue}} \approx 5\times10^{-31}\,\text{m}^2\) (for λ≈450 nm)  
    * \(L = 1000\ \text{m}\)
    
    \[
    \tau_{\text{blue}} \approx (2.5\times10^{25})(5\times10^{-31})(10^3) \approx 0.013
    \]
    
    So about **1 %** of blue photons are scattered in a 1 km slab. Over the whole atmosphere (≈ 10 km effective scattering height) the cumulative scattering reaches a few percent—enough to give the sky its characteristic brightness.
    
    ---
    
    ## TL;DR for a non‑technical audience
    
    Sunlight contains all colors. When it passes through the air, the tiny gas molecules act like tiny prisms that preferentially bounce around the **blue** part of the light. Because that blue light is redirected toward us from every direction, the sky looks blue. At sunrise and sunset the light has to travel through **even more** air, so most of the blue gets scattered away before it reaches us, leaving the reds and oranges we see at those times.
    
    --- 
    
    **Key take‑away:** The sky is blue because **Rayleigh scattering** makes the atmosphere far more effective at redirecting short‑wavelength (blue) sunlight than longer‑wavelength (red) sunlight, and our eyes see that scattered blue light coming from all over the dome above us.
    Response tijd in seconden: 12.76668
\end{minted}

\section{Ollama AsyncClient (offline) resultaten}

\begin{listing}
    \begin{minted}{python}
        import time
        t1 = time.time()
        
        # https://github.com/ollama/ollama-python?tab=readme-ov-file#async-client
        import asyncio
        from ollama import AsyncClient
        
        async def chat():
        message = {'role': 'user', 'content': 'Why is the sky blue?'}
        response = await AsyncClient().chat(model='gemma3', messages=[message])
        # print(type(response))
        # print(dir(response))
        print(response.message.content)
        t2 = time.time()
        print("\nResponse tijd in seconden: %.5f " % (t2 - t1))
        
        print("*********after async function call**********")
        
        asyncio.run(chat())
    \end{minted}
    \caption[Ollama AsyncClient Code]{Code om te communiceren met \textit{gemma3:4b}.}
    \label{lst:ollama-asyncclient-test}
\end{listing}

\textbf{Vraag:} 'Why is the sky blue?'

\textbf{Output:}

\begin{minted}{markdown}
    *********after async function call**********
    That's a fantastic question, and one that's fascinated people for centuries! The short answer is: **the sky is blue because of a phenomenon called Rayleigh scattering.** Here's a breakdown of the explanation:
    
    **1. Sunlight and Colors:**
    
    * Sunlight appears white, but it's actually made up of all the colors of the rainbow – red, orange, yellow, green, blue, indigo, and violet. You can see this when light passes through a prism and splits into a rainbow.
    
    **2. What is Rayleigh Scattering?**
    
    * **Atmospheric Particles:** The Earth’s atmosphere is filled with tiny particles like nitrogen and oxygen molecules.
    * **Scattering Light:** When sunlight enters the atmosphere, it collides with these particles. This collision causes the light to scatter in different directions.
    * **Wavelength Matters:**  Here's the key:  **Shorter wavelengths of light (blue and violet) are scattered *much* more strongly than longer wavelengths (red and orange).** Think of it like throwing a small ball (blue light) at a bumpy surface – it's more likely to bounce off in many directions. A larger ball (red light) would be less affected.
    
    **3. Why Blue Specifically?**
    
    * **Violet is scattered *even* more than blue:**  Violet light has the shortest wavelength and is scattered the most. However, our eyes are more sensitive to blue, and the sun emits less violet light, so we primarily see blue.
    
    **4. At Sunset/Sunrise:**
    
    * When the sun is low on the horizon, sunlight has to travel through *much* more of the atmosphere to reach your eyes.
    * This longer path means that almost all the blue light has been scattered away before it reaches you.
    * The remaining light is mostly the longer wavelengths – orange and red – which haven't been scattered as much.  That's why sunsets and sunrises often appear red or orange.
    
    
    
    **In short:  The sky is blue because blue light is scattered more effectively by the molecules in the air than other colors.**
    
    ---
    
    **Resources for more information:**
    
    * **NASA - Why is the sky blue?:** [https://science.nasa.gov/sky-and-telescope/optics/why-is-the-sky-blue/](https://science.nasa.gov/sky-and-telescope/optics/why-is-the-sky-blue/)  
    * **National Geographic - Why is the sky blue?:** [https://www.nationalgeographic.org/encyclopedia/sky](https://www.nationalgeographic.org/encyclopedia/sky)
    
    
    Do you want me to delve into any specific aspect of this explanation, such as:
    
    *   The role of different wavelengths?
    *   How this affects the color of clouds?
    
    Response tijd in seconden: 67.68256
\end{minted}

\section{Ollama Async Generator (offline) resultaten}

\begin{listing}
    \begin{minted}{python}
        import time
        t1 = time.time()
        
        # https://github.com/ollama/ollama-python?tab=readme-ov-file#async-client
        import asyncio
        from ollama import AsyncClient
        
        async def chat():
        message = {'role': 'user', 'content': 'Why is the sky blue?'}
        async for part in await AsyncClient().chat(model='gemma3', messages=[message], stream=True):
        print(part['message']['content'], end='', flush=True)
        
        asyncio.run(chat())
        
        t2 = time.time()
        print("\nResponse tijd in seconden: %.5f " % (t2 - t1))
    \end{minted}
    \caption[Ollama Async Generator Code]{Code om te communiceren met \textit{gemma3:4b}.}
    \label{lst:ollama-async-generator-test}
\end{listing}

\textbf{Vraag:} 'Why is the sky blue?'

\textbf{Output:}

\begin{minted}{markdown}
    Okay, let's break down why the sky is blue – it’s a fascinating phenomenon called **Rayleigh scattering**. Here's the explanation:
    
    **1. Sunlight and its Colors:**
    
    * Sunlight isn't actually white. It's made up of *all* the colors of the rainbow – red, orange, yellow, green, blue, indigo, and violet.  You can see this when light passes through a prism and separates into its different colors.
    
    **2. The Atmosphere and Scattering:**
    
    * The Earth's atmosphere is filled with tiny particles – mostly nitrogen and oxygen molecules.
    * When sunlight enters the atmosphere, it collides with these molecules. This collision causes the light to scatter in different directions.
    
    **3. Rayleigh Scattering – The Key Process:**
    
    * **Rayleigh scattering** is a specific type of scattering that’s *most effective with shorter wavelengths of light*.  Shorter wavelengths are blue and violet light.
    * Think of it like this: blue and violet light waves are much smaller and more energetic. They are more easily bounced around by the tiny air molecules.  Red and orange light have longer wavelengths and are less affected – they tend to pass straight through.
    
    **4. Why Blue, Not Violet?**
    
    * Violet light is scattered *even more* than blue. However, there are a couple of reasons we see blue instead of violet:
    * **Sunlight's Spectrum:** The sun emits slightly less violet light than blue light.
    * **Our Eyes:** Our eyes are more sensitive to blue light than violet.
    
    
    **In simpler terms:**  The air molecules scatter blue light more than other colors, so when we look up at the sky, we’re seeing the scattered blue light.
    
    **What about sunsets?**
    
    At sunset (or sunrise), the sunlight has to travel through *much* more of the atmosphere to reach our eyes.  This long path means that almost all of the blue light has been scattered away *before* it gets to us.  The longer wavelengths (reds and oranges) are less scattered and can make it through, giving us those beautiful sunset colors.
    
    
    
    **Resources for Further Learning:**
    
    * **NASA - Why is the sky blue?** [https://science.nasa.gov/sky-science/why-is-the-sky-blue/](https://science.nasa.gov/sky-science/why-is-the-sky-blue/)
    * **Wikipedia - Rayleigh scattering:** [https://en.wikipedia.org/wiki/Rayleigh_scattering](https://en.wikipedia.org/wiki/Rayleigh_scattering)
    
    
    Do you want to delve into any specific aspect of this explanation, such as:
    
    *   The math behind Rayleigh scattering?
    *   How it relates to other colors like red and orange?
    Response tijd in seconden: 69.89036
\end{minted}

\section{Ollama Img Async Generator (offline) resultaten}

\begin{listing}
    \begin{minted}{python}
        import time
        t1 = time.time()
        
        # https://github.com/ollama/ollama-python?tab=readme-ov-file#async-client
        
        from ollama import chat
        import base64
        
        MODEL1 = 'qwen3-vl:8b-instruct'
        MODEL2 = 'gemma3'
        
        img1 = r"C:\Users\peter\OneDrive\Bureaublad\Workspace\Ollama\Beksinski 01.jpg"
        img2 = r"C:\Users\peter\OneDrive\Bureaublad\Workspace\Ollama\untitled-nightmare-tree-zdzislaw-beksinski.jpg"
        
        with open(img2, "rb") as img:
        s = base64.b64encode(img.read()).decode("utf-8")
        
        stream = chat(
        model=MODEL2,
        messages=[{'role': 'user', 'content': "what's in this image?", 'images': [s]}],
        stream=True,
        )
        
        for chunk in stream:
        print(chunk['message']['content'], end='', flush=True)
        
        t2 = time.time()
        print("\nResponse tijd in seconden: %.5f " % (t2 - t1))
    \end{minted}
    \caption[Ollama Img Async Generator Code]{Code om te communiceren met \textit{gemma3:4b}.}
    \label{lst:ollama-img-async-generator-test}
\end{listing}

\textbf{Vraag:} 'what's in this image?' + afbeelding als base64 encoded tekst

\textbf{Output:}

\begin{minted}{markdown}
    Here's a breakdown of what's depicted in the image:
    
    *   **Distorted Tree:** The most striking element is a massive, deeply shadowed tree with branches that curve dramatically upward and outward, almost like a monstrous face. It dominates the scene.
    
    *   **Moon and Sky:** Above the tree is a crescent moon in a pale, hazy sky. The sky has a yellowish/golden hue.
    
    *   **Red Landscape:** Below the tree is a vast, desolate landscape covered in a thick layer of deep red foliage, like dried leaves or perhaps a blood-like substance.
    
    *   **Boat-Like Structures:** Two dark, boat-like or vaguely humanoid structures are present on the red landscape, adding to the surreal and unsettling atmosphere.
    
    **Overall Impression:** The painting evokes a sense of unease, melancholy, and the surreal. It is strongly reminiscent of the style of Zdzisław Beksiński, a Polish surrealist painter known for his dark, unsettling, and often apocalyptic landscapes.
    Response tijd in seconden: 86.98152
\end{minted}

\section{Ollama Localhost-API Img (offline) resultaten}

\begin{listing}
    \begin{minted}{python}
        import base64
        import requests
        import ollama
        
        import time
        t1 = time.time()
        
        '''
        code from: https://ollama.com/library/llava  &  https://github.com/ollama/ollama/blob/main/docs/api.md#request-5
        
        Request:
        
        curl http://localhost:11434/api/generate -d '{
            "model": "llava",
            "prompt":"What is in this picture?",
            "images": ["iVBORw0KGgoAAAANSUhEUgAA..."]
        }'
        
        Response:
        
        {
            "model": "llava",
            "created_at": "2023-11-03T15:36:02.583064Z",
            "response": "A happy cartoon character, which is cute and cheerful.",
            "done": true,
            "context": [1, 2, 3],
            "total_duration": 2938432250,
            "load_duration": 2559292,
            "prompt_eval_count": 1,
            "prompt_eval_duration": 2195557000,
            "eval_count": 44,
            "eval_duration": 736432000
        }
        
        '''
        
        MODEL = 'llava:7b'
        
        def model_installed(model) -> bool:
        models = ollama.list().models
        model_names = [m.model for m in models]
        if model in model_names:
        print("Model is installed, executing query.")
        return True
        print(f"Model not yet installed, installing {model}")
        return False
        
        
        if not model_installed(MODEL):
        print(f"Please wait while Installing {MODEL}")
        ollama.pull(MODEL)
        else:
        print(f"Encoding image and prompting {MODEL}")
        # Encode the image to a base64 string.
        with open(r"C:\Users\peter\OneDrive\Bureaublad\Workspace\Ollama\Beksinski 01.jpg", "rb") as img:
        s = base64.b64encode(img.read()).decode("utf-8")
        
        json = {
            "model": MODEL,
            "prompt": "What is in this picture?",
            "stream": False,
            "images": [s]
        }
        response = requests.post(
        "http://localhost:11434/api/generate",
        json=json
        )
        # print(response.text) # (response.status_code / response.headers / response.content / response.text / response.json())
        
        data = response.json()
        '''
        print(data.keys())
        dict_keys(['model', 'created_at', 'response', 'done', 'done_reason', 'context', 'total_duration', 'load_duration', 'prompt_eval_count', 'prompt_eval_duration', 'eval_count', 'eval_duration'])
        '''
        text_output = data.get("response", "")
        print(text_output)
        
        t2 = time.time()
        print("\nResponse tijd in seconden: %.5f " % (t2 - t1))
    \end{minted}
    \caption[Ollama API Img Code]{Code om te communiceren met \textit{llava:7b}.}
    \label{lst:ollama-localhost-api-img-test}
\end{listing}

\textbf{Vraag:} 'What is in this picture?' + afbeelding als base64 encoded tekst

\textbf{Output:}

\begin{minted}{markdown}
    Model is installed, executing query.
    Encoding image and prompting llava:7b
    The image features a large, stylized tree with a prominent canopy and multiple branches. It appears to be set within an artificial or painted environment that mimics a forest setting. In the background, there's a glowing light source emanating from between the branches of the tree, creating a somewhat magical or ethereal atmosphere. The color palette is rich with warm tones, contributing to the serene and mystical ambiance of the scene.
    
    Response tijd in seconden: 73.27115
\end{minted}

\section{Ollama + Tools (offline) resultaten}

Zie: \ref{sssec:ollamaplustools}

\section{Ollama + JavaScript (offline) resultaten}

\begin{listing}
    \begin{minted}{js}
        /**
        * setup
        * -----
        * 
        * ollama model downloaden & ollama programma installeren binnen de node omgeving
        * 
        * bash:
        * ollama pull gemma3
        * npm i ollama
        */
        
        /* js */
        import ollama from 'ollama'
        
        const STARTTIME = new Date().getTime();
        
        async function runDemo() {
            /* codefragment from: https://docs.ollama.com/quickstart#javascript */
            const response = await ollama.chat({
                model: 'gemma3',
                messages: [{ role: 'user', content: 'Why is the sky blue?' }],
            })
            console.log(response.message.content)
            /* end codefragment */
            
            const ENDTIME = new Date().getTime();
            console.log(`\nResponse tijd in seconden: ${(ENDTIME - STARTTIME) / 1000}`);
        }
        
        runDemo();
        
        /* demo opstarten:
        Bash:
        cd containing-directory
        node quickstart.js
        */
    \end{minted}
    \caption[Ollama Javascript Code]{Code om te communiceren met \textit{gemma3:4b}.}
    \label{lst:ollama-javascript-test}
\end{listing}

\textbf{Vraag:} 'Why is the sky blue?'

\textbf{Output:}

\begin{minted}{markdown}
    (node:40768) [MODULE_TYPELESS_PACKAGE_JSON] Warning: Module type of file:///C:/Users/peter/OneDrive/Bureaublad/Workspace/Ollama/js/quickstart.js is not specified and it doesn't parse as CommonJS.
    Reparsing as ES module because module syntax was detected. This incurs a performance overhead.
    To eliminate this warning, add "type": "module" to C:\Users\peter\package.json.
    (Use `node --trace-warnings ...` to show where the warning was created)
    That's a fantastic question, and it's one that's fascinated people for centuries! The short answer is: **the sky is blue because of a phenomenon called Rayleigh scattering.** Here's a more detailed explanation:
    
    **1. Sunlight and its Colors:**
    
    * Sunlight appears white, but it’s actually made up of *all* the colors of the rainbow (red, orange, yellow, green, blue, indigo, and violet). You can see this when you see a rainbow – those colors are separated by raindrops.
    
    **2. What is Rayleigh Scattering?**
    
    * **Light and Air Molecules:** The Earth's atmosphere is full of tiny air molecules (mostly nitrogen and oxygen) and particles.
    * **Scattering:** When sunlight enters the atmosphere, it collides with these air molecules. This collision causes the light to scatter in different directions.     
    * **Wavelength Matters:** Here's the crucial part: **shorter wavelengths of light (blue and violet) are scattered *much* more strongly than longer wavelengths (red and orange).**  This is because the shorter wavelengths interact more powerfully with the small air molecules.
    
    **3. Why Blue, Not Violet?**
    
    * Violet light has the shortest wavelength and is scattered *even more* than blue. However, there are a couple of reasons why we see a blue sky instead of a violet one:
    * **Sun's Emission:** The sun actually emits slightly less violet light than blue light.
    * **Our Eyes are Less Sensitive to Violet:** Our eyes are more sensitive to blue light than violet.
    
    
    **Think of it like this:** Imagine throwing a small ball (blue light) and a large ball (red light) at a bumpy surface (the atmosphere). The small ball will bounce around in many directions, while the larger ball will mostly continue in its original path.
    
    **4. Sunsets & Sunrises – Why Red and Orange?**
    
    At sunrise and sunset, the sun is lower on the horizon. This means the sunlight has to travel through *much* more of the atmosphere to reach our eyes.  Because the blue light has been scattered away along its longer path, the remaining light that reaches us is predominantly the longer wavelengths - red and orange.
    
    
    
    **Resources for Further Learning:**
    
    * **NASA - Why is the sky blue?:** [https://science.nasa.gov/sky-science/why-is-the-sky-blue/](https://science.nasa.gov/sky-science/why-is-the-sky-blue/)
    * **National Geographic - Why is the Sky Blue?:** [https://www.nationalgeographic.org/encyclopedia/sky](https://www.nationalgeographic.org/encyclopedia/sky)
    
    
    Do you want me to delve deeper into any specific aspect of this explanation, such as:
    
    *   The physics behind Rayleigh scattering in more detail?
    *   How this affects seeing stars?
    *   The role of aerosols (pollution) in the sky’s color?
    
    Response tijd in seconden: 60.498
\end{minted}

\section{LangChain}

\subsection{LangChain Eerste Code Voorbeeld}

\begin{listing}
    \begin{minted}{python}
        from langchain_ibm import WatsonxLLM
        
        model_id = "meta-llama/llama-3-405b-instruct"
        
        parameters = {
            GenParams.MAX_NEW_TOKENS: 256,  # this controls the maximum number of tokens in the generated output
            GenParams.TEMPERATURE: 0.5, # this randomness or creativity of the model's responses
        }
        
        url = "https://us-south.ml.cloud.ibm.com"
        project_id = "skills-network"
        
        llm = WatsonxLLM(
            model_id=model_id,
            url=url,
            project_id=project_id,
            params=parameters
        )
        
        template = """Tell me a {adjective} joke about {content}.
        """
        prompt = PromptTemplate.from_template(template)
        # type(prompt) -> langchain_core.prompts.prompt.PromptTemplate
        
        prompt.format(adjective="funny", content="chickens")
        # -> 'Tell me a funny joke about chickens.\n'
        
        from langchain_core.runnables import RunnableLambda
        
        # Define a function to ensure proper formatting
        def format_prompt(variables):
            return prompt.format(**variables)
        
        # Create the chain with explicit formatting
        joke_chain = (
            RunnableLambda(format_prompt)
            | llm 
            | StrOutputParser()
        )
        
        # Run the chain
        response = joke_chain.invoke({"adjective": "funny", "content": "chickens"})
        print(response)
        '''
        output:
        Why did the chicken go to the doctor?
        Because it had fowl breath!
        I hope that one made you cluck with laughter!
        '''
        
        response = joke_chain.invoke({"adjective": "sad", "content": "fish"})
        print(response)
        '''
        output:
        Why did the fish go to the party?
        Because he heard it was a "reel" good time, but he got hooked on the idea of having friends and ended up feeling crabby when nobody showed up. 🐟😔
        I hope that made a splash of sadness in your heart! 💦😢
        '''
    \end{minted}
    \caption[Eerste voorbeeld van code met LangChain]{Een eerste code voorbeeld van LangChain, \autocite{IBM2026}}
    \label{lst:langchaineerstecodevoorbeeld}
\end{listing}

Dit codefragment gebruikt ook een LLM van IBM, zij hebben een eigen custom LLM-wrapper ontwikkeld met IBM-Granite.

\subsection{LangChain Verschillende use case code voorbeelden}

\subsubsection{Code Generatie:}
\begin{figure}[H]
    \centering
    \includegraphics[width=1.0\textwidth]{LCEL_code-generation.png}
    \caption[Code generatie AI agent met LangChain]{\label{fig:lcelcodegenlangchainagent}Code generatie AI agent met LangChain en LCEL, \autocite{IBM2026}}
\end{figure}
\underline{Output:} \\

\begin{minted}{markdown}
    SELECT c.name, c.email 
    FROM customers c 
    JOIN purchases p ON c.customer_id = p.customer_id 
    WHERE p.purchase_date >= NOW() - INTERVAL 30 DAY;
    
    
    Note: The database system being used is MySQL. The query uses the JOIN clause to combine rows from the 'customers' and 'purchases' tables based on matching 'customer_id' values. The WHERE clause filters the results to include only purchases made in the last 30 days. The NOW() function returns the current date and time, and the INTERVAL keyword is used to subtract 30 days from this value. 
    
    
    If you want to write a Python program to generate this SQL query, you can use the following code:
    
    ```python
    def generate_sql_query(table_name, days):
        # Define the base query
        query = f"SELECT c.name, c.email FROM {table_name} c JOIN purchases p ON c.customer_id = p.customer_id WHERE "
        
        # Add the date filter condition
        query += f"p.purchase_date >= NOW() - INTERVAL {days} DAY;"
        
        return query
    
    # Generate the SQL query
    sql_query = generate_sql_query('customers', 30)
    print(sql_query
\end{minted}

\subsubsection{Tekst Samenvatting:}
\begin{figure}[H]
    \centering
    \includegraphics[width=1.0\textwidth]{LCEL_tekst-samenvatting.png}
    \caption[tekst samenvatting AI agent met LangChain]{\label{fig:lceltextsummarylangchainagent}Tekst samenvatting AI agent met LangChain en LCEL, \autocite{IBM2026}}
\end{figure}
\underline{Output:} \\
The rapid advancement of technology in the 21st century has transformed various industries, including healthcare, education, and transportation, through innovations like AI, machine learning, and the Internet of Things.


\subsubsection{Tekst Classificatie:}
\begin{figure}[H]
    \centering
    \includegraphics[width=1.0\textwidth]{LCEL_tekst-classificatie.png}
    \caption[tekst samenvatting AI agent met LangChain]{\label{fig:lceltextclassificationlangchainagent}Tekst classificatie AI agent met LangChain en LCEL, \autocite{IBM2026}}
\end{figure}
\underline{Output:} \\
Music \\
Reasoning Skill: \\
The reasoning skill for this question is Emotion Recognition And Sentiment Analysis. The sentence expresses a positive sentiment towards the concert, using words like "exhilarating" and "outstanding". The correct classification into the Music category requires recognizing the emotional tone of the sentence and understanding that it is describing a musical event. This involves analyzing the language used and identifying the sentiment expressed, which is a key aspect of Emotion Recognition And Sentiment Analysis.
